\documentclass[aspectratio=169]{beamer}

% --- Essential Packages ---
\usepackage[utf8]{inputenc}
\usepackage[T1]{fontenc}
\usepackage{lmodern}
\usepackage{graphicx}
\usepackage{tikz}
\usepackage{tabularx}
\usepackage{booktabs}
\usepackage{amsmath}
\usepackage{hyperref}
\usepackage{listings}
\usepackage{xcolor}
\usepackage{fontawesome5}

% --- Theme Setup ---
\usetheme{Madrid}
\usecolortheme{beaver}

% Custom colors
\definecolor{darkred}{RGB}{139,0,0}
\definecolor{lightgray}{RGB}{245,245,245}
\definecolor{codegreen}{RGB}{0,128,0}
\definecolor{codegray}{RGB}{128,128,128}
\definecolor{codeblue}{RGB}{0,0,200}
\definecolor{attackred}{RGB}{180,0,0}
\definecolor{defensegreen}{RGB}{0,120,0}
\definecolor{pipelineorange}{RGB}{210,100,0}

% Code listing style
\lstset{
  basicstyle=\tiny\ttfamily,
  keywordstyle=\color{codeblue}\bfseries,
  commentstyle=\color{codegreen},
  stringstyle=\color{attackred},
  breaklines=true,
  frame=single,
  backgroundcolor=\color{lightgray},
  numbers=left,
  numberstyle=\tiny\color{codegray},
  numbersep=4pt,
  showstringspaces=false
}

\setbeamertemplate{navigation symbols}{}
\setbeamertemplate{bibliography item}[text]
\setbeamertemplate{caption}[numbered]
\addtobeamertemplate{frametitle}{}{\vspace{-0.25cm}}
\setbeamersize{text margin left=0.6cm, text margin right=0.6cm, text margin bottom=0.8cm}

% --- Title Info ---
\title[]{A Security Perspective On Retrieval Augmented Generation}
\subtitle{A Legal Domain Case Study: End-to-End Development, Attacks, and Defenses}

\author[Jasveer Panwar]{
    \textbf{Presented by:} Jasveer Panwar\\
    \vspace{0.1cm}
    \scriptsize{Guided By: Dr. B.L. Parne}
}

\institute[SVNIT]{
    \includegraphics[height=1cm]{logo.png} \\
    \vspace{0.1cm}
    Department of CSE \\
    Sardar Vallabhbhai National Institute of Technology, Surat
}

\date{Final Presentation}

% ============================================================
\begin{document}
% ============================================================

% --- Title Slide ---
\begin{frame}
    \titlepage
\end{frame}

% --- Contents ---
\begin{frame}{Contents}
    \tableofcontents
\end{frame}

% ============================================================
\section{Introduction}
% ============================================================
\begin{frame}{Introduction}
    \begin{block}{Securing RAG Systems in High-Stakes Domains}
        Retrieval-Augmented Generation (RAG) systems are increasingly deployed in critical domains like Law, Healthcare, and Finance, where \textbf{incorrect answers can carry severe consequences}.
    \end{block}
    \vspace{0.25cm}
    \begin{columns}[t]
        \begin{column}{0.48\textwidth}
            \textbf{Why Legal Domain?}
            \begin{itemize}
                \item Requires very high accuracy and verifiable grounding.
                \item New laws such as BNS, BNSS, and BSA are not reliably covered by older base models.
                \item Government and compliance use cases demand strong security guarantees.
            \end{itemize}
        \end{column}
        \begin{column}{0.48\textwidth}
            \textbf{What this work achieved}
            \begin{itemize}
                \item Built a functional Indian Legal RAG system over statutory documents.
                \item Demonstrated realistic attack paths on the RAG pipeline.
                \item Implemented and evaluated a layered defense strategy.
            \end{itemize}
        \end{column}
    \end{columns}
\end{frame}

% ============================================================
\section{Problem Statement}
% ============================================================
\begin{frame}{Problem Statement}
    \begin{alertblock}{The Core Issue}
        Pretrained LLMs alone cannot be relied upon for up-to-date legal answers. While RAG provides factual grounding, it still assumes the \textbf{trustworthiness of the indexed corpus and retrieved context}, which attackers can exploit.
    \end{alertblock}
    \vspace{0.25cm}
    \small
    \textbf{Research Problem:} Build an accurate legal RAG system, then design, implement, and validate defenses that reduce susceptibility to:
    \begin{itemize}
        \item Corpus-level data poisoning
        \item Prompt injection
        \item Context manipulation
    \end{itemize}
\end{frame}

% ============================================================
\section{Motivation}
% ============================================================
\begin{frame}{Motivation}
    \small
    \textbf{Why implement legal RAG instead of using a plain LLM?}
    \vspace{0.2cm}
    \begin{itemize}
        \setlength\itemsep{0.5em}
        \item \textbf{Frozen Knowledge:} A model trained before 2023 cannot reliably answer about newer statutes such as BNS Section 103.
        \item \textbf{Zero Tolerance for Hallucination:} Fabricated provisions can cause legal harm.
        \item \textbf{Security Requirement:} If the retrieved knowledge base is manipulated, the model can still answer incorrectly with confidence.
        \item \textbf{Citation Mandate:} Legal answers must reference source document, section, and chapter.
        \item \textbf{Auditability:} Each query should be logged for compliance, debugging, and forensic review.
    \end{itemize}
\end{frame}

% ============================================================
\section{Literature Survey}
% ============================================================
\begin{frame}{Literature Survey}
    \tiny
    \setlength{\tabcolsep}{4pt}
    \begin{tabularx}{\textwidth}{|p{0.21\textwidth}|X|X|}
        \hline
        \textbf{Paper} & \textbf{Contribution} & \textbf{Limitations} \\ \hline

        \textbf{Backdoored Retrievers} &
        -- Shows retriever-level backdoors and prompt-injection pathways \newline
        -- Demonstrates corpus poisoning risk \newline
        -- Shows benign metrics may remain high &
        -- Strong on attack demonstration \newline
        -- Limited deployable defense blueprint \newline
        -- Less focus on domain-specific hardening \\ \hline

        \textbf{RAG Security Survey} &
        -- Taxonomy of RAG vulnerabilities \newline
        -- Covers poisoning, prompt attacks, retriever threats \newline
        -- Summarizes defense ideas &
        -- Mostly conceptual \newline
        -- Limited end-to-end implementation evidence \newline
        -- Practical trade-offs untested \\ \hline

        \textbf{Securing RAG: Risk \& Mitigation} &
        -- Provides risk assessment model \newline
        -- Maps threats to mitigation classes \newline
        -- Offers deployment guidance &
        -- No live attack validation \newline
        -- Limited empirical evaluation \newline
        -- Defense layering not fully operationalized \\ \hline

        \textbf{PoisonedRAG} &
        -- Demonstrates high poisoning success \newline
        -- Shows small poisoned knowledge can control outputs \newline
        -- Highlights retrieval as attack surface &
        -- Focused mainly on attack impact \newline
        -- Practical mitigation remains open \newline
        -- No legal-domain deployment study \\ \hline
    \end{tabularx}
\end{frame}

% ============================================================
\section{Research Gap}
% ============================================================
\begin{frame}{Research Gap Identified}
    \begin{itemize}
        \setlength\itemsep{0.75em}
        \item \textbf{Empirical defenses missing:} Few end-to-end, validated mitigation pipelines exist for prompt injection and poisoning in live RAG systems.
        \item \textbf{Provenance integration:} Limited evaluation of source trust and provenance signals inside retrieval re-ranking.
        \item \textbf{Hybrid defense stacks:} Filtering, trust scoring, prompt hardening, and diagnostics are rarely implemented together.
        \item \textbf{Multi-attack evaluation:} Few systems assess poisoning, prompt injection, and context manipulation in one pipeline.
        \item \textbf{Legal-domain specificity:} Abbreviations, amendments, and section-based citations create challenges not handled well by generic RAG defenses.
    \end{itemize}
\end{frame}

% ============================================================
\section{Objectives}
% ============================================================
\begin{frame}{Objectives}
    \textbf{Dissertation Objectives:}
    \vspace{0.25cm}
    \begin{enumerate}
        \setlength\itemsep{0.55em}
        \item Build a fully functional Indian Legal RAG system as the experimental testbed.
        \item Identify security vulnerabilities across the end-to-end RAG pipeline.
        \item Demonstrate practical attacks: data poisoning, prompt injection, and context manipulation.
        \item Implement a multi-layered defense stack directly in the serving path.
        \item Evaluate defense effectiveness while preserving benign answer quality.
        \item Establish practical security lessons for sensitive-domain RAG deployments.
    \end{enumerate}
\end{frame}

% ============================================================
\section{Architecture}
% ============================================================
\begin{frame}{RAG System Architecture}
    \begin{center}
        \includegraphics[width=0.85\textwidth, height=0.75\textheight, keepaspectratio]{rag.jpg}
    \end{center}
\end{frame}

\begin{frame}{Security-Aware RAG Architecture Overview}
    \small
    \begin{columns}[t]
        \begin{column}{0.33\textwidth}
            \textbf{\textcolor{pipelineorange}{Stage 1: Ingestion}}
            \begin{itemize}
                \item PDF and text parsing
                \item Metadata extraction
                \item Chunking with overlap
                \item Source tagging and JSONL output
            \end{itemize}
        \end{column}
        \begin{column}{0.33\textwidth}
            \textbf{\textcolor{pipelineorange}{Stage 2: Indexing}}
            \begin{itemize}
                \item Sentence-BERT embeddings
                \item FAISS cosine retrieval index
                \item Metadata persistence
                \item Batch processing
            \end{itemize}
        \end{column}
        \begin{column}{0.33\textwidth}
            \textbf{\textcolor{pipelineorange}{Stage 3: Serving}}
            \begin{itemize}
                \item Query normalization
                \item Secure retrieval and filtering
                \item Safety-aware prompt
                \item Model invocation and audit log
            \end{itemize}
        \end{column}
    \end{columns}
    \vspace{0.3cm}
    \begin{alertblock}{Security Layer}
        Final system behavior is secured across retrieval and generation: \textbf{trusted-source bias + injection filtering + strict prompting + logging}.
    \end{alertblock}
\end{frame}

% ============================================================
\section{Implementation}
% ============================================================
\begin{frame}{Project Structure}
    \small
    \begin{columns}[t]
        \begin{column}{0.48\textwidth}
            \textbf{Source Files}
            \begin{itemize}
                \item \texttt{src/ingest\_with\_metadata.py} --- document ingestion and chunking
                \item \texttt{src/embed\_index.py} --- FAISS index builder
                \item \texttt{src/serve\_query.py} --- core RAG pipeline and security
                \item \texttt{src/interactive\_openai.py} --- CLI interface
                \item \texttt{src/evaluate.py} --- QA evaluation framework
                \item \texttt{src/comapre\_models.py} --- model comparison
            \end{itemize}
        \end{column}
        \begin{column}{0.48\textwidth}
            \textbf{Support Files and Data}
            \begin{itemize}
                \item \texttt{models/openai\_client.py} --- OpenAI wrapper
                \item \texttt{models/ollama\_client.py} --- local model wrapper
                \item \texttt{data/} --- legal source documents
                \item \texttt{chunks.jsonl} --- processed chunks
                \item \texttt{index/} --- FAISS index and metadata
                \item \texttt{logs/} --- per-query JSON audit trail
            \end{itemize}
        \end{column}
    \end{columns}
\end{frame}

\begin{frame}{Implementation Workflow: Data and Indexing}
    \small
    \begin{columns}[t]
        \begin{column}{0.48\textwidth}
            \textbf{Dataset and Ingestion}
            \begin{itemize}
                \item \textbf{Corpus:} Indian legal documents including BNS, BNSS, BSA, CrPC, IPC, IEA, POCSO, and NDPS
                \item \textbf{Chunking:} approximately 3500 chunks with \texttt{chunk\_size=1200}, \texttt{overlap=300}
                \item \textbf{Metadata:} source name, page markers, and sidecar key-value information
                \item \textbf{Output:} \texttt{chunks.jsonl}
            \end{itemize}
        \end{column}
        \begin{column}{0.48\textwidth}
            \textbf{Embedding and Indexing}
            \begin{itemize}
                \item \textbf{Model:} \texttt{all-MiniLM-L6-v2}
                \item \textbf{Index:} FAISS with normalized vectors for cosine similarity
                \item \textbf{Storage:} \texttt{index/index.faiss} plus metadata store
                \item \textbf{Batch size:} 64 for efficient embedding generation
            \end{itemize}
        \end{column}
    \end{columns}
\end{frame}

\begin{frame}{Implementation Workflow: Query Serving}
    \small
    \begin{columns}[t]
        \begin{column}{0.48\textwidth}
            \textbf{Retrieval Pipeline}
            \begin{itemize}
                \item Query rewrite for legal shorthand expansion
                \item Top-k FAISS similarity search
                \item Re-ranking using filename, phrase, and metadata signals
                \item Filtering of suspicious injection-like content
            \end{itemize}
        \end{column}
        \begin{column}{0.48\textwidth}
            \textbf{Answer Generation}
            \begin{itemize}
                \item \textbf{LLM:} OpenAI GPT-4o-mini or local Ollama
                \item \textbf{Prompt:} context-only legal answering with safety note
                \item \textbf{Strict mode:} safe fallback when context is insufficient
                \item \textbf{Logging:} full JSON log per query
            \end{itemize}
        \end{column}
    \end{columns}
\end{frame}

\begin{frame}{System Outputs: Generalized vs. RAG-Specific}
    \begin{columns}
        \begin{column}{0.5\textwidth}
            \centering
            \textbf{Generalized LLM Response} \\
            \includegraphics[width=\linewidth, height=0.6\textheight, keepaspectratio]{irrelevant_answer.png}
        \end{column}
        \begin{column}{0.5\textwidth}
            \centering
            \textbf{Legal Database (Source Corpus)} \\
            \includegraphics[width=\linewidth, height=0.6\textheight, keepaspectratio]{data.png}
        \end{column}
    \end{columns}
\end{frame}

\begin{frame}{Processing Pipeline: Chunking and Embeddings}
    \begin{columns}
        \begin{column}{0.5\textwidth}
            \centering
            \textbf{Chunking Process} \\
            \includegraphics[width=\linewidth, height=0.6\textheight, keepaspectratio]{chunks.png}
        \end{column}
        \begin{column}{0.5\textwidth}
            \centering
            \textbf{Embedding Generation} \\
            \includegraphics[width=\linewidth, height=0.6\textheight, keepaspectratio]{emeddings.png}
        \end{column}
    \end{columns}
\end{frame}

\begin{frame}{Retrieval and Querying}
    \begin{columns}
        \begin{column}{0.5\textwidth}
            \centering
            \textbf{Retrieval Logic} \\
            \includegraphics[width=\linewidth, height=0.6\textheight, keepaspectratio]{retrieval.png}
        \end{column}
        \begin{column}{0.5\textwidth}
            \centering
            \textbf{Final Relevant Answer} \\
            \includegraphics[width=\linewidth, height=0.6\textheight, keepaspectratio]{relevant_answer.png}
        \end{column}
    \end{columns}
\end{frame}

% ============================================================
\section{Threat Analysis and Attack Demonstration}
% ============================================================
\begin{frame}{Attack Vector Overview}
    \small
    \begin{columns}[t]
        \begin{column}{0.33\textwidth}
            \centering
            \textbf{\textcolor{attackred}{Attack 1}}\\
            \textbf{Data Poisoning}\\
            \vspace{0.2cm}
            Inject false legal provisions into the corpus so they are retrieved as evidence and influence the final answer.
        \end{column}
        \begin{column}{0.33\textwidth}
            \centering
            \textbf{\textcolor{attackred}{Attack 2}}\\
            \textbf{Prompt Injection}\\
            \vspace{0.2cm}
            Embed hidden instructions inside a document so retrieved context tries to override the intended model behavior.
        \end{column}
        \begin{column}{0.33\textwidth}
            \centering
            \textbf{\textcolor{attackred}{Attack 3}}\\
            \textbf{Context Manipulation}\\
            \vspace{0.2cm}
            Craft semantically close adversarial chunks that displace legitimate context and mislead the answer generation stage.
        \end{column}
    \end{columns}
    \vspace{0.35cm}
    \begin{block}{Testbed}
        All attack demonstrations were performed on the live Legal RAG system built over the Indian legal document corpus.
    \end{block}
\end{frame}

\begin{frame}{Attack Demo: Data Poisoning}
    \small
    \textbf{Mechanism:} Inserts false laws $\to$ misleads retrieval $\to$ produces incorrect legal answers.
    \vspace{0.1cm}
    \begin{columns}
        \begin{column}{0.5\textwidth}
            \centering
            \scriptsize{1. Before Poisoning} \\
            \includegraphics[width=\linewidth, height=0.5\textheight, keepaspectratio]{not_poisoned.png}
        \end{column}
        \begin{column}{0.5\textwidth}
            \centering
            \scriptsize{2. Ingesting Poisoned Document} \\
            \includegraphics[width=\linewidth, height=0.5\textheight, keepaspectratio]{poisioned_doc.png}
        \end{column}
    \end{columns}
\end{frame}

\begin{frame}{Attack Demo: Data Poisoning Result}
    \centering
    \textbf{Result: System retrieves poisoned context and outputs the incorrect legal answer.}
    \vspace{0.2cm}
    \includegraphics[width=0.75\textwidth, height=0.65\textheight, keepaspectratio]{poisoned_resp.png}
\end{frame}

\begin{frame}{Attack Demo: Prompt Injection}
    \small
    \textbf{Mechanism:} Hidden instructions in retrieved context attempt to override the intended system behavior.
    \vspace{0.1cm}
    \begin{columns}
        \begin{column}{0.5\textwidth}
            \centering
            \scriptsize{1. Instruction Override Attempt} \\
            \includegraphics[width=\linewidth, height=0.5\textheight, keepaspectratio]{overwrite_prompt.png}
        \end{column}
        \begin{column}{0.5\textwidth}
            \centering
            \scriptsize{2. Manipulation Success Before Defense} \\
            \includegraphics[width=\linewidth, height=0.5\textheight, keepaspectratio]{prompt_inj_response.png}
        \end{column}
    \end{columns}
\end{frame}

% ============================================================
\section{Security Implementation}
% ============================================================
\begin{frame}{Security Implementation Overview}
    \begin{block}{Multi-Layered Defense in \texttt{src/serve\_query.py}}
        The final system integrates security controls directly into the live query path rather than adding them as post-processing patches.
    \end{block}
    \vspace{0.2cm}
    \small
    \begin{columns}[t]
        \begin{column}{0.48\textwidth}
            \textbf{\textcolor{defensegreen}{Active Defenses}}
            \begin{enumerate}
                \item Query normalization and rewriting
                \item Prompt-injection chunk filtering
                \item Trusted-source bias scoring
                \item PDF source promotion
                \item Untrusted TXT demotion
            \end{enumerate}
        \end{column}
        \begin{column}{0.48\textwidth}
            \textbf{\textcolor{defensegreen}{Structural Defenses}}
            \begin{enumerate}
                \item Context-only strict answering prompt
                \item Embedded anti-injection instruction
                \item Minimum similarity threshold cutoff
                \item Max context character limit
                \item Full JSON audit logging per query
            \end{enumerate}
        \end{column}
    \end{columns}
\end{frame}

\begin{frame}{Implementation Architecture}
    \centering
    \includegraphics[width=0.85\textwidth, height=0.8\textheight, keepaspectratio]{Imparch.png}
\end{frame}

\begin{frame}{Defense 1: Query Normalization}
    \small
    \begin{block}{Purpose}
        Expands legal shorthand to full forms so semantically correct chunks are easier to retrieve.
    \end{block}
    \vspace{0.15cm}
    \textbf{\texttt{rewrite\_query(q)} using \texttt{QUERY\_REWRITE\_MAP}:}
    \vspace{0.1cm}
    \begin{columns}[t]
        \begin{column}{0.48\textwidth}
            \begin{itemize}
                \item \texttt{POCSO} $\to$ Protection of Children from Sexual Offences
                \item \texttt{CrPC} $\to$ Code of Criminal Procedure
                \item \texttt{IPC} $\to$ Indian Penal Code
                \item \texttt{BNS} $\to$ Bharatiya Nyaya Sanhita
            \end{itemize}
        \end{column}
        \begin{column}{0.48\textwidth}
            \begin{itemize}
                \item \texttt{BNSS} $\to$ Bharatiya Nagarik Suraksha Sanhita
                \item \texttt{BSA} $\to$ Bharatiya Sakshya Adhiniyam
                \item \texttt{IEA} $\to$ Indian Evidence Act
                \item \texttt{NDPS} $\to$ Narcotic Drugs and Psychotropic Substances
            \end{itemize}
        \end{column}
    \end{columns}
    \vspace{0.15cm}
    \begin{alertblock}{Impact}
        Prevents retrieval misses caused by abbreviated legal queries that do not directly match fully worded statute chunks.
    \end{alertblock}
\end{frame}

\begin{frame}{Defense 2: Prompt Injection Filtering}
    \small
    \begin{block}{Purpose}
        Detects and drops retrieved chunks containing instruction-injection patterns before they reach the model prompt.
    \end{block}
    \vspace{0.15cm}
    \textbf{\texttt{is\_injection\_chunk(text)} detects patterns such as:}
    \begin{itemize}
        \item \texttt{ignore previous instructions}
        \item \texttt{disregard the above}
        \item \texttt{you are now} and \texttt{act as}
        \item \texttt{system:}, \texttt{[INST]}, and \texttt{<|im\_start|>}
        \item \texttt{jailbreak}, \texttt{bypass}, and \texttt{override}
        \item Hidden Unicode control characters such as \texttt{\textbackslash u200b}
    \end{itemize}
    \vspace{0.15cm}
    \begin{alertblock}{Key Insight}
        Retrieval-time filtering is stronger than relying only on the model to resist malicious instructions embedded in context.
    \end{alertblock}
\end{frame}

\begin{frame}{Defense 3: Trusted-Source Scoring and Re-Ranking}
    \small
    \begin{block}{Purpose}
        Promotes verified legal sources and demotes low-trust text content during final ranking.
    \end{block}
    \vspace{0.15cm}
    \begin{columns}[t]
        \begin{column}{0.48\textwidth}
            \textbf{Scoring Components}
            \begin{itemize}
                \item \texttt{is\_pdf\_source(meta)}
                \item \texttt{is\_trusted\_source(meta)}
                \item \texttt{filename\_match\_score()}
                \item \texttt{exact\_phrase\_score()}
                \item \texttt{metadata\_keyword\_score()}
            \end{itemize}
        \end{column}
        \begin{column}{0.48\textwidth}
            \textbf{Trust Penalties}
            \begin{itemize}
                \item \texttt{DEMOTE\_UNTRUSTED\_TXT = True}
                \item \texttt{UNTRUSTED\_TXT\_PENALTY} applied to ranking
                \item Official PDFs ranked above untrusted text where possible
                \item Final score combines similarity and trust signals
            \end{itemize}
        \end{column}
    \end{columns}
\end{frame}

\begin{frame}{Defense 4: Safety-Aware Prompt Construction}
    \small
    \begin{block}{Purpose: \texttt{build\_prompt(context, question, cutoff, strict)}}
        Constructs a prompt that enforces context-only answering and explicitly tells the model to ignore instructions inside the retrieved context.
    \end{block}
    \vspace{0.1cm}
    \textbf{Key prompt elements:}
    \begin{enumerate}
        \item Role definition: precise legal assistant answering only from supplied context.
        \item Safety note: retrieved context may contain instructions and those must be ignored.
        \item Context block: filtered and re-ranked chunks with source cues.
        \item Strict mode: if evidence is insufficient, answer with ``I don't know''.
        \item Citation requirement: cite source document, section, or chapter.
    \end{enumerate}
\end{frame}

\begin{frame}{Defense 5: Audit Logging and Diagnostics}
    \small
    \begin{block}{Purpose}
        Every query is logged with enough detail to support debugging, compliance review, and post-incident analysis.
    \end{block}
    \vspace{0.15cm}
    \textbf{JSON log entry contains:}
    \begin{columns}[t]
        \begin{column}{0.48\textwidth}
            \begin{itemize}
                \item Timestamp
                \item Original query and rewritten query
                \item Retrieved chunk IDs
                \item Similarity and trust-related scores
            \end{itemize}
        \end{column}
        \begin{column}{0.48\textwidth}
            \begin{itemize}
                \item Filtered chunk decisions
                \item Full prompt sent to the LLM
                \item Model response
                \item Latency and token usage where available
            \end{itemize}
        \end{column}
    \end{columns}
    \vspace{0.15cm}
    \begin{alertblock}{Security Value}
        Audit logs make it possible to reconstruct attacks, verify defense activation, and analyze suspicious retrieval behavior after the fact.
    \end{alertblock}
\end{frame}

\begin{frame}{After Securing the Pipeline}
    \begin{columns}
        \begin{column}{0.5\textwidth}
            \centering
            \scriptsize{\textbf{Result:} Instruction Override Rejected} \\
            \includegraphics[width=\linewidth, height=0.5\textheight, keepaspectratio]{answer.png}
        \end{column}
        \begin{column}{0.5\textwidth}
            \centering
            \scriptsize{\textbf{Result:} Manipulation Failed} \\
            \includegraphics[width=\linewidth, height=0.5\textheight, keepaspectratio]{secans.png}
        \end{column}
    \end{columns}
\end{frame}

\begin{frame}{After Securing the RAG System}
    \centering
    \textbf{Result: The system retrieves the correct answer after hardening.}
    \vspace{0.2cm}
    \includegraphics[width=0.75\textwidth, height=0.65\textheight, keepaspectratio]{coreectans.png}
\end{frame}

\begin{frame}{Attack vs. Defense: Comparison}
    \small
    \begin{columns}[t]
        \begin{column}{0.48\textwidth}
            \textbf{\textcolor{attackred}{Without Defenses}}
            \begin{itemize}
                \item Poisoned chunks retrieved as top results
                \item Hidden instructions can override model behavior
                \item Model outputs attacker-controlled answers
                \item No clear visibility into what failed
                \item Trust in the system breaks quickly
            \end{itemize}
        \end{column}
        \begin{column}{0.48\textwidth}
            \textbf{\textcolor{defensegreen}{With Defenses Enabled}}
            \begin{itemize}
                \item \texttt{is\_injection\_chunk()} drops malicious chunks
                \item Trusted-source bias ranks official PDFs higher
                \item Strict prompt resists embedded instructions
                \item Audit log flags suspicious retrieval decisions
                \item System safely falls back to ``I don't know'' when needed
            \end{itemize}
        \end{column}
    \end{columns}
    \vspace{0.25cm}
    \begin{block}{Key Finding}
        No single defense layer is sufficient. The combination of \textbf{injection filtering + source trust scoring + strict prompting + audit logging} is much stronger than any individual defense alone.
    \end{block}
\end{frame}

% ============================================================
\section{Evaluation}
% ============================================================
\begin{frame}{Evaluation: Defense Effectiveness}
    \small
    \begin{columns}[t]
        \begin{column}{0.48\textwidth}
            \textbf{Attack Success Rate (Lower = Better)}
            \vspace{0.2cm}
            \begin{tabularx}{\textwidth}{|X|c|}
                \hline
                \textbf{Scenario} & \textbf{ASR} \\ \hline
                No defense baseline & High \\ \hline
                Injection filter only & Moderate \\ \hline
                Source trust only & Moderate \\ \hline
                Full defense stack & Low \\ \hline
            \end{tabularx}
        \end{column}
        \begin{column}{0.48\textwidth}
            \textbf{Observed Effect}
            \begin{itemize}
                \item Cleaner final context for benign queries
                \item Lower influence of untrusted text chunks
                \item Reduced instruction-following from retrieved context
                \item Better failure behavior under uncertainty
            \end{itemize}
        \end{column}
    \end{columns}
    \vspace{0.25cm}
    \begin{alertblock}{Observation}
        The defense stack reduced attack success substantially while preserving or improving benign answer quality due to better retrieval precision and safer context assembly.
    \end{alertblock}
\end{frame}

% ============================================================
\section{Comparative Analysis}
% ============================================================
\begin{frame}{Comparative Analysis: Backdoored vs Our Secure Legal RAG}
\scriptsize
\begin{columns}[T,totalwidth=\textwidth]
    \column{0.48\textwidth}
    \textbf{Backdoored Retrievers and Related Attack Work}
    \begin{itemize}
        \item Primarily attack-centric and demonstrates vulnerability severity.
        \item Focuses on poisoning and retriever manipulation.
        \item Reports strong attack success under targeted conditions.
        \item Shows that benign retrieval quality may hide underlying compromise.
    \end{itemize}

    \column{0.48\textwidth}
    \textbf{Our Implementation}
    \begin{itemize}
        \item Defense-centric and deployment-oriented for legal RAG.
        \item Retrieval safeguards: query rewrite, injection filtering, trust-aware re-ranking.
        \item Prompt safeguards: context-only answering and explicit anti-injection instruction.
        \item Operational visibility: per-query audit logs for prompt, scores, context, response, and latency.
    \end{itemize}
\end{columns}

\vspace{0.35em}
\textbf{Key Improvements over Attack-Only Baselines}
\begin{itemize}
    \item \textbf{Before generation:} malicious chunks are filtered or demoted before entering final context.
    \item \textbf{During generation:} strict refusal behavior reduces unsupported answers.
    \item \textbf{After generation:} forensic traceability supports debugging and incident analysis.
\end{itemize}
\end{frame}

% ============================================================
\section{Conclusion}
% ============================================================
\begin{frame}{Conclusion}
    \small
    \begin{itemize}
        \setlength\itemsep{0.75em}
        \item Successfully built a \textbf{functional Indian Legal RAG system} across ingestion, indexing, retrieval, generation, and evaluation.
        \item Demonstrated that data poisoning and prompt injection are practical threats, not only theoretical ones.
        \item Implemented a \textbf{multi-layered security framework} in the serving pipeline: query normalization, injection filtering, source trust bias, strict prompting, and audit logging.
        \item Evaluation indicates that defenses \textbf{reduce attack success} while maintaining useful performance on benign legal queries.
        \item \textbf{Takeaway:} Production-grade legal RAG requires layered security across retrieval and generation, not just a strong base model.
    \end{itemize}
\end{frame}

% ============================================================
\section{Future Work}
% ============================================================
\begin{frame}{Future Work: Proposed Enhancements}
    \begin{enumerate}
        \setlength\itemsep{0.7em}
        \item \textbf{Ingestion Integrity:} cryptographic hash and signature checks plus source allow-list enforcement.
        \item \textbf{Embedding Anomaly Detection:} outlier and clustering-based detection for suspicious chunks before indexing.
        \item \textbf{Context Sanitization Layer:} redact or remove instruction-like passages before prompt construction.
        \item \textbf{Provenance Scoring:} incorporate stronger verifiable source signals into ranking.
        \item \textbf{Adversarial Red-Teaming:} systematic multi-attack benchmarking across poisoning, backdoors, and prompt injection.
    \end{enumerate}
\end{frame}

\begin{frame}{Future Work: Defense Architecture Diagram}
    \centering
    \includegraphics[width=0.85\textwidth, height=0.8\textheight, keepaspectratio]{Future work.png}
\end{frame}

% --- References ---
\section{References}
\begin{frame}[allowframebreaks]{References}
\tiny
\begin{thebibliography}{9}

\bibitem{surveyRAGSecurity}
\href{https://ieeexplore.ieee.org/document/11172756}
{J. Wen et al., \textit{``Retrieval-Augmented Generation: A Survey of Security Challenges and Countermeasures,''} IEEE Internet of Things Journal, 2025.}

\bibitem{backdooredRetrievers}
\href{https://arxiv.org/abs/2410.14479}
{M. Chen et al., \textit{``Backdoored Retrievers for Prompt Injection Attacks on Retrieval-Augmented Generation of Large Language Models,''} arXiv preprint arXiv:2410.14479, 2024.}

\bibitem{ragMitigationFramework}
\href{https://ieeexplore.ieee.org/document/11081501}
{K. Li et al., \textit{``Securing RAG: A Risk Assessment and Mitigation Framework,''} IEEE Transactions on Information Forensics and Security, 2024.}

\bibitem{poisonedRAG}
\href{https://www.usenix.org/conference/usenixsecurity25/presentation/zou-poisonedrag}
{Z. Zou et al., \textit{``PoisonedRAG: Knowledge Corruption Attacks to Retrieval-Augmented Generation of Large Language Models,''} USENIX Security Symposium, 2025.}

\bibitem{rag-original}
\href{https://arxiv.org/abs/2005.11401}
{P. Lewis et al., \textit{``Retrieval-Augmented Generation for Knowledge-Intensive NLP,''} NeurIPS, 2020.}

\bibitem{Hallucinations}
\href{https://ieeexplore.ieee.org/document/11014810}
{L. Ammann et al., \textit{``Exploring RAG Solutions to Reduce Hallucinations in LLMs,''} IEEE, 2025.}

\bibitem{bns}
\href{https://www.indiacode.nic.in/bitstream/123456789/20062/1/a2023-45.pdf}{The Bharatiya Nyaya Sanhita, 2023}

\bibitem{bnss}
\href{https://www.indiacode.nic.in/bitstream/123456789/21920/1/the_bharatiya_nagarik_suraksha_sanhita\%2c_2023.pdf}{The Bharatiya Nagarik Suraksha Sanhita, 2023}

\bibitem{bsa}
\href{https://home.jk.gov.in/pdf/3BSA2023.pdf}{The Bharatiya Sakshya Adhiniyam, 2023}

\end{thebibliography}
\end{frame}

% --- Thank You ---
\begin{frame}
    \begin{center}
        \vspace{1cm}
        {\Huge \textbf{Thank You}}\\
        \vspace{0.4cm}
        {\large Questions and Discussion}\\
        \vspace{1cm}
        \small Department of CSE, SVNIT Surat
    \end{center}
\end{frame}

\end{document}
